\section{Discussion}
\label{sec:discussion}

\subsection{Scope of the Online-Learning Claim}

A recurring temptation in self-adaptive systems is to let strong
scenarios launder weak evidence. Our pre-submission review caught the
beginning of exactly this. The clean RAN experiment --- the principal
experiment of the paper --- triggers \emph{zero} evolution updates,
so it cannot, by itself, demonstrate that online learning improves
real-world outcomes; it demonstrates the opposite-flavored result
that the trigger does not misfire. The adaptation claim therefore
rests on the surge experiment alone, and we have scoped it that way:
under a $+35\%$ step with a causal information budget, self-evolution
halves post-surge violations across seeds. We stress what the surge
experiment is and is not. It is a controlled, repeatable regime-shift
probe on top of real demand data --- the standard methodology for
evaluating drift response when the held-out window is (fortunately)
drift-free. It is not evidence that UCRA would already have adapted
to the real test window's natural conditions; on that window there
was nothing to adapt to, which is itself the deployment-happy case.
Future work should accumulate natural-drift evidence: longer windows,
multiple network segments (the RAN dataset ships three), and
production A/B shadow comparisons where ethics and contracts permit.

\subsection{What the Leakage Audit Does and Does Not Guarantee}

The causal protocol of Section~\ref{sec:stage5} guarantees that the
learner's training labels at update time $t$ are a subset of the
labels observable at $t$. Three caveats bound the guarantee. First,
it covers \emph{label}-availability leakage; it does not
automatically detect other leakage channels (e.g., preprocessing
statistics fitted on the full window), which our audit handles by
construction (scalers fitted on train segments only) but a different
pipeline might not. Second, the surge experiment's \emph{input
windows} are built from the drifted series, which is correct (the
deployed system would observe drifted inputs) but means the
experiment measures adaptation under the assumption that drift is
observable in the inputs --- a stealth drift that alters outcomes
without touching observable demand could evade Trigger B, though
Trigger A (violations) would still catch it in operation. Third,
causality is necessary but not sufficient for good adaptation: the
protocol ensures the evaluation is honest, not that the fine-tuning
hyperparameters are optimal; our sensitivity grid suggests the
update behavior is robust (identical trigger points across seeds and
across fine-tune hyperparameter settings in the released grid), but
we present the residual post-surge violations as the honest cost of
gentle adaptation rather than a tuned optimum.

\subsection{Why a Simple Transform Can Beat a Cheating Oracle}

The most counterintuitive row in Table~\ref{tab:main} is the oracle's
16.3\% violations. The mechanism is worth spelling out because it
generalizes. A rolling historical quantile reacts to bursts
\emph{after} they enter its window and forgets lulls as the window
advances; on heavy-tailed 15-minute demand, its conditional violation
probability sits well above the nominal $1 - \tau$. UCRA's
reservation, by contrast, is anchored to a \emph{learned} conditional
distribution that maps contextual structure (hour of day, weekend,
recent volatility) to the quantile, and its kappa buffer adds slack
exactly where the learned band is wide. The lesson for practitioners
is that nominal quantile levels are not violation guarantees unless
the quantiles are conditional, calibrated, and buffered --- and the
kappa dial is precisely the operator's handle for converting
``approximately calibrated'' into ``operationally zero.''

\subsection{Limitations and Threats to Validity}

\emph{External validity.} All RAN-side results come from one
18.5-day segment of one network (dataset01 of three available);
three seeds quantify training variance, not network variance. The
CTTC workload is simulation-derived with its own traffic
assumptions. Neither dataset exercises multi-domain (radio +
transport + core) reservation, and the capacity model
($1.3\times$ peak) is an operator rule, not a measured envelope.

\emph{Construct validity.} Violation is binary per slot; real SLAs
grade severity (duration, depth, affected users), and a policy could
game the metric with tiny-but-frequent violations. Our logged
per-slot artifacts enable severity-aware re-analysis, but the
headlines here are binary.

\emph{Internal validity.} The leakage audit covers the three
channels we found; the regression tests pin them. We cannot prove
the absence of unknown channels --- only that the final pipeline's
information flows were reviewed end to end, and that the audit
methodology is released with the code for re-review.

\emph{Statistical scope.} Three seeds is a pragmatic budget for a
52k-parameter model trained in minutes, adequate to expose the
$\pm 3$--6 point spreads we report, but marginal for tail claims
(e.g., the exact knee of the $\kappa$ frontier). The released sweep
artifacts include per-seed curves so readers can apply their own
inference.

\emph{Baseline scope.} We compare against reservation policies that
share UCRA's information budget. We do not compare against
reinforcement-learning slice controllers; doing so would strengthen
the auditability argument empirically but requires an exploration
budget on live infrastructure that is outside this study's ethics
and scope. The train-only $\Phi$ and oracle quantile baselines
bracket what forecast-free and forecast-omniscient policies can do,
which is the comparison UCRA's design thesis requires.

\emph{Deployment economics.} On the evaluated segment, static peak
sizing idles 30.5\% of capacity against UCRA's 10.7\% --- a 19.8-point
reduction in waste, i.e., UCRA commits $\approx$65\% less idle
capacity for the identical zero-violation outcome. At 15-minute
granularity this compounds to roughly 2{,}470 slot-capacity
equivalents freed per 1{,}000 slots of operation; at sector scale and
backhaul tariffs this is the difference between buying and
re-timing capacity. The remaining 10.7\% waste is not slop: it is
the priced insurance premium of Proposition~\ref{prop:bound}
(the buffer above the q90 base), it is stable across seeds
($13.3\% \pm 2.1$), and it can be traded toward utilization with
Table~\ref{tab:kappa} in hand rather than discovered in an outage
post-mortem.

\subsection{Generalization Beyond Radio Demand}

The UCRA loop assumes nothing specific about radio access: it needs
a univariate demand stream at a fixed slot, a capacity envelope, and
feedback about whether the reservation held. The same five stages
therefore transfer, with re-parameterization, to transport-bandwidth
reservations between data centers, compute reservations for
cloud-native network functions (where the CTTC-style snapshot
evaluation applies naturally), energy procurement for base-site
batteries under time-of-use tariffs, and standby-capacity sizing for
network operations centers. What does not transfer automatically is
the forecast's seasonality structure (24-hour windows encode a
diurnal regime; weekly data would want $K$ re-tuned) and the drift
trigger thresholds, which are policy choices tied to each domain's
violation economics. The framework's claim is the contract between
uncertainty and commitment --- calibration in, auditable reservation
out, causal adaptation in between --- not the radio domain
itself.

\subsection{Deployment Considerations}

The deployed loop is deliberately light: one LSTM forward pass
(52k parameters, sub-millisecond on CPU), linear transform
arithmetic, a 512-window buffer, and at most a few fine-tune events
per day under drift. The operational levers map to organizational
roles: $\kappa$ and $\tau$ are policy (owned by SLA/finance),
trigger thresholds are operations (owned by NOC), and the forecast
model is engineering (retrained offline, patched online by Stage~5).
The all-zero-violation operating point's $\approx$28\% utilization
gain over peak sizing translates directly into capacity that can be
resold or deferred --- on the evaluated segment, roughly 28\% of
three weeks of sector capacity --- which is the economic argument
for the entire pipeline.

\begin{figure}[!t]
\centering
\diagramplaceholder{4.2cm}{deployment-context}
{Suggested content: a deployment topology diagram --- the UCRA loop
running beside the network's PM counter pipeline: PM counters
$\rightarrow$ ingestion $\rightarrow$ UCRA controller (Stages 1--5)
$\rightarrow$ reservation commands to the slice manager / OSS;
drift alerts and $\kappa$ policy console for the NOC; and the offline
retraining path. Mark which components are online (per-slot) versus
offline (per-seed).}
\caption{Deployment context of the UCRA loop.}
\label{fig:deployment}
\end{figure}
